\chapter*{Tài liệu tham khảo}
\addcontentsline{toc}{chapter}{Tài liệu tham khảo}
\label{ch:tham-khao}

Báo cáo này mô tả một hệ thống đang chạy được, nên nguồn của nó là chính mã nguồn và các bản ghi
quyết định đi kèm, chứ không phải tài liệu bên ngoài. Mọi con số và mọi luật nêu trong báo cáo
đều có một chỗ tương ứng trong kho dưới đây.

\begin{description}[itemsep=1.1em,parsep=0.35em,leftmargin=0pt]

  \item[Mã nguồn và tài liệu kỹ thuật của dự án]\mbox{}\\
  \url{\linkrepo}
  \medskip

  Kho chứa năm service đã mô tả ở chương~\ref{ch:kien-truc-va-du-lieu}, cùng ba nhóm tài liệu
  mà báo cáo này rút gọn từ đó:

  \begin{itemize}
    \item \texttt{docs/decisions/} --- các bản ghi quyết định nghiệp vụ. Mỗi quyết định ở đó ghi
          cả phần \emph{đã cân nhắc những gì} và \emph{nơi nào thi hành nó}, hai phần mà báo cáo
          này không chở.
    \item \texttt{docs/overview/} --- mô tả nghiệp vụ, kiến trúc và mô hình dữ liệu. Nguồn của
          chương~\ref{ch:luong-nghiep-vu}, \ref{ch:use-case} và~\ref{ch:kien-truc-va-du-lieu}.
    \item \texttt{docs/diagrams/} --- tệp gốc của mọi sơ đồ trong báo cáo. Các hình ở đây được
          \emph{xuất ra} từ chúng chứ không vẽ lại, và một phép kiểm tự động trong kho sẽ báo đỏ
          nếu một sơ đồ đổi mà hình chưa xuất lại.
  \end{itemize}

  \item[Các luật được canh bằng máy]\mbox{}\\
  \texttt{tools/check\_contract.py} và \texttt{AGENTS.md} trong kho trên.
  \medskip

\end{description}
